%\includepdf[
%scale=1,
%pages=1,
%pagecommand={\section{Anhang}\subsection{Planung und Konzeption des Proof of Concept}\label{anhang:poc}\thispagestyle{plain}}
%]{./anhang/Planung_und_Konzeption_des_PoC.pdf}

% So könnte man voll PDF als Anhang nutzen, ist aber unschön.

% !TEX root = ../master.tex

\appendix
\section{Anhang}
\label{sec:anhang}

\subsection{Verzeichnis der portierten Lehrnotebooks}
\label{app:notebook_liste}

Die folgende Tabelle listet alle Jupyter-Notebooks auf, die im Rahmen dieser Arbeit selektiert und von Python nach TypeScript portiert wurden. Die Gliederung folgt der Kapitelstruktur der Originalvorlesung \enquote{Theoretische Informatik III}.

\begin{table}[h]
	\caption{Übersicht der übersetzten Notebooks im Repository}
	\centering
	\renewcommand{\arraystretch}{1.2}
	\begin{tabularx}{\textwidth}{@{}lX@{}}
		\toprule
		\textbf{Original-Kapitel} & \textbf{Notebook-Dateiname} \\
		\midrule
		\multicolumn{2}{@{}l}{\textbf{Chapter 03}} \\
		& \texttt{01-Ply-Scanning-Example.ipynb} \\
		& \texttt{02-Regexp-Tutorial.ipynb} \\
		& \texttt{03-Exam-Evaluation.ipynb} \\
		& \texttt{04-Html2Text.ipynb} \\
		\addlinespace
		\multicolumn{2}{@{}l}{\textbf{Chapter 04 \& 05}} \\
		& \texttt{01-NFA-2-DFA.ipynb} \\
		& \texttt{02-Test-NFA-2-DFA.ipynb} \\
		& \texttt{03-Regexp-2-NFA.ipynb} \\
		& \texttt{04-Test-Regexp-2-NFA.ipynb} \\
		& \texttt{05-DFA-2-RegExp.ipynb} \\
		& \texttt{06-Test-DFA-2-RegExp.ipynb} \\
		& \texttt{07-Minimize.ipynb} \\
		& \texttt{08-Test-Minimize.ipynb} \\
		& \texttt{09-Equivalence.ipynb} \\
		& \texttt{10-Test-Equivalence.ipynb} \\
		\addlinespace
		\multicolumn{2}{@{}l}{\textbf{Chapter 06}} \\
		& \texttt{01-Top-Down-Parser.ipynb} \\
		& \texttt{02-EBNF-Parser.ipynb} \\
		\addlinespace
		\multicolumn{2}{@{}l}{\textbf{Chapter 07}} \\
		& \texttt{01-Calculator.ipynb} \\
		& \texttt{02-Differentiator.ipynb} \\
		& \texttt{03-Interpreter.ipynb} \\
		\addlinespace
		\multicolumn{2}{@{}l}{\textbf{Chapter 10}} \\
		& \texttt{01-Conflicts.ipynb} \\
		& \texttt{02-Conflicts-Resolved.ipynb} \\
		\bottomrule
	\end{tabularx}
	\vspace{0,2 cm}
	\raggedright \footnotesize \textbf{Quelle:} Eigene Darstellung.
	\label{tab:notebook_list}
\end{table}